% An "animation" here is not motion. Each step prints as ANOTHER PAGE
% of the PDF, showing a little more. A frame with three steps becomes
% three pages. That is why an animated deck has more pages than slides.

% --- Bullets, one at a time ---------------------------------------
\begin{frame}{Built up}
  \begin{steps}
    \item Appears on page 1
    \item Appears on page 2
    \item Appears on page 3
  \end{steps}
\end{frame}

% --- Break anywhere ------------------------------------------------
\begin{frame}{Two halves}
  The setup.
  \stepafter
  The punchline.
\end{frame}

% --- Precise control -----------------------------------------------
\showfrom{2}{Appears from page 2 on; space held before it.}
\showonly{2}{Appears only on page 2; no space held.}
\emphat{2}{Highlighted on page 2, plain otherwise.}
\swap{2}{Before}{After}

% --- Swapping a figure in place ------------------------------------
\begin{frame}{Same spot, two figures}
  \showonly{1}{\slidefigure{before.pdf}}%
  \showonly{2}{\slidefigure{after.pdf}}%
\end{frame}

% Beamer's own <1->, \pause, \only and \alt still work. Use whichever
% reads better. Keep animation for places where the ORDER matters --
% every step is a page someone has to sit through.

% --- A figure that stays while the text builds ---------------------
% The classic trap: put the overlays INSIDE a pane, not around the
% columns, or the whole split appears only on the last page.
\begin{frame}{Figure holds, text builds}
  \begin{figuretext}[0.55]
    \begin{steps}
      \item First point
      \item Second point
    \end{steps}
    \nextpane
    \slidefigure{diagram.pdf}   % no overlay spec: present throughout
  \end{figuretext}
\end{frame}

% Printing an animated deck? See snippets/handout.tex -- one word
% collapses every step back to a single page per slide.
